\noindent\textbf{Einheit 3 · Prozentwert berechnen}

\erg{16}{a) 3 von 10 Kästchen grau \quad b) 3 von 4 Kästchen grau \quad c) ein halbes von 10 Kästchen grau \quad d) etwa sechs Zehntel des Streifens grau}
\erg{17}{a) $40\,€$ \quad b) $15\,$kg \quad c) $20\,$m \quad d) $18\,$l}
\erg{18}{a) $30\,€$ \quad b) $20\,$kg \quad c) $7\,$m \quad d) $23\,€$ \quad e) $21\,€$ \quad f) $12\,€$ \quad g) $49\,$kg \quad h) $3{,}00\,€$ \quad i) $60\,$kg}
\erg{19}{a) $9\,€$ \quad b) $36\,€$ \quad c) $45{,}60\,€$ \quad d) $48\,€$ (der Restpreis wäre $272\,€$)}
\erg{20}{a) Lena hat $40\,\%$ als $0{,}04$ statt als $0{,}4$ geschrieben; richtig ist $0{,}4 \cdot 60\,€ = 24\,€$ \quad b) $63\,€$}
\erg{21}{a) $1\,\%$ ist der hundertste Teil des Ganzen; von diesem einen Prozent kann man jeden Prozentsatz durch Malnehmen bekommen \quad b) ja – beide Male rechnet man $0{,}25 \cdot 80$ bzw. $0{,}8 \cdot 25$, und die Reihenfolge der Faktoren ändert das Produkt nicht ($20\,€$)}
\erg{22}{a) $10\,800$ \quad b) $9\,000$ \quad c) Altstadt; der Prozentsatz zählt vom eigenen Ganzen, und Altstadt hat ein viel größeres Ganzes}
